% ============================================================
%  C: Como Programar -- 6ª Edição | Deitel & Deitel
%  Capítulo 1 -- Introdução aos Computadores, à Internet e à Web
%  EXERCÍCIOS (1.3 a 1.9)
%  Estilo didático e incremental -- no espírito de Paul Deitel
% ============================================================
\documentclass[12pt,a4paper]{article}

\usepackage[utf8]{inputenc}
\usepackage[T1]{fontenc}
\usepackage[brazil]{babel}
\usepackage{geometry}
\geometry{top=2.5cm, bottom=2.5cm, left=3cm, right=3cm}
\usepackage{parskip}
\setlength{\parskip}{6pt}
\setlength{\parindent}{0pt}


\usepackage{xcolor}
\usepackage{tcolorbox}
\tcbuselibrary{skins, breakable, theorems}
\usepackage{enumitem}
\usepackage{listings}
\usepackage{fancyhdr}
\usepackage{titlesec}
\usepackage{hyperref}
\usepackage{booktabs}
\usepackage{array}

% Paleta Deitel
\definecolor{deitelblue}{RGB}{0, 70, 127}
\definecolor{deitelgray}{RGB}{240, 242, 245}
\definecolor{deitelgreen}{RGB}{0, 110, 60}
\definecolor{deitellorange}{RGB}{200, 90, 0}
\definecolor{deitellightblue}{RGB}{220, 235, 250}
\definecolor{deitelred}{RGB}{160, 20, 20}
\definecolor{verdadeiro}{RGB}{0, 100, 50}
\definecolor{falso}{RGB}{160, 20, 20}

% Caixas
\newtcolorbox{boaprática}[1][]{colback=deitellightblue,colframe=deitelblue,
  fonttitle=\bfseries\small,title={Boa prática de programação},breakable,#1}
\newtcolorbox{erroco}[1][]{colback=red!5,colframe=deitelred,
  fonttitle=\bfseries\small,title={Erro comum de programação},breakable,#1}
\newtcolorbox{portab}[1][]{colback=green!5,colframe=deitelgreen,
  fonttitle=\bfseries\small,title={Dica de portabilidade},breakable,#1}
\newtcolorbox{respostabox}[1][]{colback=deitelgray,colframe=deitelblue!60,
  fonttitle=\bfseries\small\color{deitelblue},title={Resposta detalhada},breakable,#1}
\newtcolorbox{enunciadoBox}[1][]{colback=white,colframe=deitelblue,
  fonttitle=\bfseries,title={#1},breakable}
\newtcolorbox{verdadeiroBox}{colback=green!5,colframe=verdadeiro,
  fonttitle=\bfseries\color{verdadeiro},title={VERDADEIRO},breakable}
\newtcolorbox{falsoBox}{colback=red!5,colframe=falso,
  fonttitle=\bfseries\color{falso},title={FALSO},breakable}

\setlist[enumerate,1]{label=\textbf{\alph*)},leftmargin=2em}
\setlist[itemize]{leftmargin=2em}

\lstset{
  basicstyle=\ttfamily\small,
  backgroundcolor=\color{deitelgray},
  frame=single,framerule=0.4pt,
  rulecolor=\color{deitelblue!40},
  breaklines=true,
  showstringspaces=false,
  keywordstyle=\color{deitelblue}\bfseries,
  commentstyle=\color{deitelgreen}\itshape,
  stringstyle=\color{deitelred},
  language=C,
  numbers=none,
  xleftmargin=10pt,xrightmargin=5pt,
}

\pagestyle{fancy}
\fancyhf{}
\fancyhead[L]{\small\color{deitelblue}\textbf{C: Como Programar -- 6ª ed.\ $|$ Deitel \& Deitel}}
\fancyhead[R]{\small\color{deitelblue}Cap.\ 1 -- Exercícios}
\fancyfoot[C]{\small\thepage}
\renewcommand{\headrulewidth}{0.4pt}

\titleformat{\section}[block]{\large\bfseries\color{deitelblue}}{\thesection.}{0.5em}{}[\titlerule]
\titleformat{\subsection}[block]{\normalsize\bfseries\color{deitelblue!80}}{\thesubsection.}{0.5em}{}

\hypersetup{colorlinks=true,linkcolor=deitelblue,urlcolor=deitelblue}

% ============================================================
\begin{document}

% ---- Capa ----
\begin{titlepage}
  \centering
  \vspace*{2cm}
  {\color{deitelblue}\rule{\linewidth}{2pt}}\\[0.4cm]
  {\huge\bfseries\color{deitelblue} C: Como Programar}\\[0.2cm]
  {\large\color{deitelblue} 6ª Edição $\cdot$ Paul Deitel \& Harvey Deitel}\\[0.2cm]
  {\color{deitelblue}\rule{\linewidth}{2pt}}\\[1cm]
  {\LARGE\bfseries Capítulo 1}\\[0.3cm]
  {\Large\color{deitelblue!80} Introdução aos Computadores, à Internet e à Web}\\[0.8cm]
  \begin{tcolorbox}[colback=deitellightblue,colframe=deitelblue,width=0.8\linewidth]
    \centering
    {\large\bfseries Exercícios (1.3 -- 1.9)}\\[0.2cm]
    {\normalsize Resolvidos passo a passo, com explicações conceituais,\\
    no estilo incremental Deitel}
  \end{tcolorbox}
  \vfill
  {\small Conteúdo restrito ao Capítulo 1 $\cdot$ Abordagem incremental}
\end{titlepage}

\tableofcontents
\newpage

% ============================================================
\section{Exercício 1.3 -- Hardware ou Software?}
% ============================================================

\begin{enunciadoBox}[Enunciado]
Categorize os seguintes itens como \textbf{hardware} ou \textbf{software}:
\begin{enumerate}
  \item CPU
  \item compilador C++
  \item ALU
  \item pré-processador C++
  \item unidade de entrada
  \item um programa editor
\end{enumerate}
\end{enunciadoBox}

\bigskip

\begin{respostabox}

\textbf{Fundamentação conceitual:}

Antes de categorizar cada item, recordemos as definições fundamentais da Seção~1.2:

\begin{itemize}
  \item \textbf{Hardware:} Os diversos \textit{dispositivos físicos} que compõem um sistema
  de computador -- tudo que pode ser tocado: circuitos, chips, teclados, discos, etc.
  \item \textbf{Software:} Os \textit{programas} executados em um computador -- as
  instruções que controlam o hardware e o fazem realizar tarefas úteis.
\end{itemize}

\medskip
\begin{center}
\renewcommand{\arraystretch}{1.5}
\begin{tabular}{>{\bfseries}c l l p{6cm}}
  \toprule
  \textbf{Item} & \textbf{Elemento} & \textbf{Categoria} & \textbf{Justificativa} \\
  \midrule
  a & CPU & \textcolor{deitelblue}{\bfseries Hardware} &
    A \textit{Unidade Central de Processamento} é um componente físico (chip
    eletrônico) que coordena todas as operações do computador. É a ``seção
    administrativa'' do hardware. \\
  \midrule
  b & Compilador C++ & \textcolor{deitellorange}{\bfseries Software} &
    Um compilador é um \textit{programa} (conjunto de instruções) que traduz código
    em linguagem de alto nível (C++) para linguagem de máquina. Reside em disco e
    é executado na CPU -- portanto, é software. \\
  \midrule
  c & ALU & \textcolor{deitelblue}{\bfseries Hardware} &
    A \textit{Unidade Lógica e Aritmética} é a ``seção de processamento'' física do
    computador, responsável por realizar cálculos e comparações. Nos computadores
    modernos, está integrada fisicamente à CPU no mesmo chip. \\
  \midrule
  d & Pré-processador C++ & \textcolor{deitellorange}{\bfseries Software} &
    O pré-processador é um \textit{programa} que processa diretivas como
    \texttt{\#include} e \texttt{\#define} antes da compilação. É software, pois
    consiste em instruções executadas pelo computador. \\
  \midrule
  e & Unidade de entrada & \textcolor{deitelblue}{\bfseries Hardware} &
    A unidade de entrada é a ``seção receptora'' física do computador -- abrange
    dispositivos como teclado, mouse, scanner. São componentes físicos, portanto,
    hardware. \\
  \midrule
  f & Programa editor & \textcolor{deitellorange}{\bfseries Software} &
    Um programa editor (como \texttt{vi}, \texttt{emacs} ou o editor do Visual Studio)
    é um \textit{programa} -- uma coleção de instruções de software -- usado para
    criar e modificar o código-fonte. \\
  \bottomrule
\end{tabular}
\end{center}

\medskip
\textbf{Resumo:} Hardware $\rightarrow$ a), c), e) \qquad Software $\rightarrow$ b), d), f)

\end{respostabox}

\begin{boaprática}
  A distinção hardware/software é fundamental em programação. Todo programa que
  você escreverá em C é \textit{software}, mas seu funcionamento depende inteiramente
  do \textit{hardware} subjacente. Entender essa relação o tornará um programador
  mais consciente sobre desempenho, portabilidade e limitações do sistema.
\end{boaprática}

\newpage

% ============================================================
\section{Exercício 1.4 -- Linguagem Independente vs.\ Dependente de Máquina}
% ============================================================

\begin{enunciadoBox}[Enunciado]
Por que você escreveria um programa em uma linguagem independente de máquina em vez
de em uma linguagem dependente da máquina? Por que uma linguagem dependente da máquina
poderia ser mais apropriada para a escrita de certos tipos de programas?
\end{enunciadoBox}

\bigskip

\begin{respostabox}

\textbf{Parte 1: Por que usar linguagem independente de máquina?}

As \textbf{linguagens de alto nível}, como C, são consideradas \textit{independentes de
máquina} porque um mesmo programa pode ser compilado e executado em diferentes tipos de
computador (com diferentes CPUs e sistemas operacionais), geralmente com poucas ou nenhuma
modificação.

As principais razões para escolher uma linguagem independente de máquina são:

\begin{enumerate}
  \item \textbf{Portabilidade:} Um programa escrito em C pode ser compilado para rodar em
  um PC com Windows, em um servidor Linux, em um Mac ou em um microcontrolador -- o mesmo
  código-fonte funciona em múltiplas plataformas. Como afirma a Dica de
  Portabilidade~1.1 do livro: \textit{``os aplicativos escritos em C podem ser executados
  com pouca ou nenhuma modificação em uma grande variedade de sistemas de computação
  diferentes''}.

  \item \textbf{Facilidade de escrita e leitura:} As linguagens de alto nível usam palavras
  próximas ao inglês e notação matemática. Um comando como
  \texttt{valorTotal = salario + horaExtra;} é imediatamente compreensível para um
  humano, ao contrário de uma sequência binária como \texttt{+1300042774}.

  \item \textbf{Produtividade:} Uma única instrução de alto nível pode equivaler a dezenas
  de instruções em linguagem de máquina ou simbólica. Isso reduz drasticamente o tempo
  de desenvolvimento de programas.

  \item \textbf{Manutenibilidade:} Programas em linguagens de alto nível são muito mais
  fáceis de ler, depurar, modificar e manter ao longo do tempo.

  \item \textbf{Reutilização:} É possível aproveitar a biblioteca-padrão de C e outras
  bibliotecas já escritas, sem se preocupar com os detalhes do hardware específico.
\end{enumerate}

\bigskip
\textbf{Parte 2: Quando uma linguagem dependente de máquina pode ser mais apropriada?}

As \textbf{linguagens dependentes de máquina}, como a linguagem de máquina e a linguagem
simbólica (assembly), apesar de difíceis de usar, oferecem vantagens em situações específicas:

\begin{enumerate}
  \item \textbf{Desempenho crítico:} O programador tem controle absoluto sobre cada
  instrução executada pelo processador. Em aplicações onde cada nanossegundo importa
  (como drivers de dispositivos, kernels de sistemas operacionais e código de tempo real),
  o assembly pode gerar código mais rápido que qualquer compilador.

  \item \textbf{Acesso direto ao hardware:} Em programação de sistemas embarcados
  (microcontroladores, dispositivos de IoT), às vezes é necessário manipular registradores,
  portas de I/O e interrupções diretamente -- operações que somente linguagens dependentes
  de máquina permitem com precisão absoluta.

  \item \textbf{Tamanho mínimo de código:} Em sistemas com memória extremamente limitada,
  o código assembly tende a ser menor que o equivalente compilado.

  \item \textbf{Inicialização de hardware:} O código de inicialização (``bootloader'')
  de um computador muitas vezes precisa ser escrito em assembly, antes que qualquer
  ambiente de linguagem de alto nível esteja disponível.
\end{enumerate}

\medskip
\textbf{Conclusão:} Para a grande maioria das aplicações, a linguagem de alto nível (como
C) é a escolha ideal, pois combina portabilidade, produtividade e legibilidade. O assembly
fica reservado para as raras situações em que o controle direto do hardware ou o desempenho
máximo são requisitos inegociáveis.

\end{respostabox}

\begin{portab}
  Curiosamente, C ocupa uma posição única: é uma linguagem de alto nível com capacidades
  de baixo nível. Por meio de ponteiros e outros recursos (estudados no Capítulo~7),
  C permite acesso próximo ao hardware sem abrir mão da portabilidade -- razão pela
  qual ainda é a linguagem predominante em sistemas operacionais e firmware.
\end{portab}

\newpage

% ============================================================
\section{Exercício 1.5 -- Preenchimento de Lacunas: Unidades do Computador}
% ============================================================

\begin{enunciadoBox}[Enunciado]
Preencha os espaços em cada uma das seguintes sentenças:
\begin{enumerate}
  \item Qual unidade lógica do computador recebe informações de fora para serem usadas por ele? \underline{\hspace{4cm}}.
  \item O processo de instruir o computador para resolver problemas específicos é chamado de \underline{\hspace{4cm}}.
  \item Que tipo de linguagem de computador usa abreviações em inglês para as instruções em linguagem de máquina? \underline{\hspace{4cm}}.
  \item Qual unidade lógica do computador envia informações que já foram processadas por ele para diversos dispositivos? \underline{\hspace{4cm}}.
  \item Quais unidades lógicas do computador retêm informações? \underline{\hspace{4cm}}.
  \item Qual unidade lógica do computador realiza cálculos? \underline{\hspace{4cm}}.
  \item Qual unidade lógica do computador toma decisões lógicas? \underline{\hspace{4cm}}.
  \item O nível de linguagem de computador mais conveniente para a elaboração rápida e fácil de programas é \underline{\hspace{4cm}}.
  \item A única linguagem que um computador entende com clareza é a chamada \underline{\hspace{4cm}} desse computador.
  \item Qual unidade lógica do computador coordena as atividades de todas as outras unidades lógicas? \underline{\hspace{4cm}}.
\end{enumerate}
\end{enunciadoBox}

\bigskip

\begin{respostabox}

\begin{enumerate}[label=\textbf{\alph*)}]

  \item \textbf{Unidade de entrada}

  A \textit{unidade de entrada} é a seção ``receptora'' do computador (Seção~1.3).
  Ela capta informações do mundo exterior -- teclado, mouse, microfone, scanner,
  sensores, dados da Internet -- e as disponibiliza para as demais unidades processarem.
  Sem ela, o computador não teria como receber dados ou instruções.

  \medskip
  \item \textbf{Programação}

  Programação é o processo pelo qual os programadores criam conjuntos ordenados de
  instruções (programas) que instruem o computador a resolver problemas específicos.
  Todo o conteúdo deste livro é dedicado a ensiná-lo a programar em C.

  \medskip
  \item \textbf{Linguagem simbólica} (também chamada de \textit{assembly})

  Da Seção~1.6: as linguagens simbólicas substituem as sequências numéricas da
  linguagem de máquina por abreviações mnemônicas próximas ao inglês, como
  \texttt{load}, \texttt{add}, \texttt{store}. Programas chamados
  \textit{assemblers} (montadores) traduzem esse código para linguagem de máquina.

  \medskip
  \item \textbf{Unidade de saída}

  A \textit{unidade de saída} é a seção de ``expedição'' do computador. Ela pega
  as informações já processadas e as envia para dispositivos externos -- monitor,
  impressora, alto-falantes, arquivos em disco, conexões de rede, etc.

  \medskip
  \item \textbf{Unidade de memória} e \textbf{unidade de armazenamento secundário}

  \textit{Duas} unidades retêm informações:
  \begin{itemize}
    \item A \textbf{unidade de memória} (RAM) armazena dados temporariamente enquanto
    o programa é executado. É volátil -- perde seu conteúdo quando o computador é
    desligado.
    \item A \textbf{unidade de armazenamento secundário} (disco rígido, SSD, etc.)
    armazena dados de forma persistente e por longos períodos.
  \end{itemize}

  \medskip
  \item \textbf{Unidade lógica e aritmética (ALU)}

  A ALU é a seção de ``processamento'' do computador. Ela executa as quatro operações
  aritméticas básicas (adição, subtração, multiplicação e divisão), bem como operações
  lógicas (comparações entre valores). Nos processadores modernos, a ALU está integrada
  dentro da CPU.

  \medskip
  \item \textbf{Unidade lógica e aritmética (ALU)}

  A mesma ALU que realiza cálculos também é responsável por \textit{tomar decisões
  lógicas} -- como comparar se dois valores são iguais, se um é maior que outro, etc.
  Essas decisões são a base de todas as estruturas de controle que você aprenderá em C
  (\texttt{if}, \texttt{while}, \texttt{for}, etc.).

  \medskip
  \item \textbf{Linguagem de alto nível}

  As linguagens de alto nível (como C, C++, Java) são as mais convenientes para o
  programador, pois usam sintaxe próxima ao inglês e permitem expressar operações
  complexas em poucas instruções. Uma única instrução de alto nível pode corresponder
  a dezenas de instruções de máquina.

  \medskip
  \item \textbf{Linguagem de máquina}

  A linguagem de máquina é o ``idioma nativo'' de cada processador específico. É a
  única linguagem que a CPU entende diretamente, sem necessidade de tradução. Consiste
  em sequências de números binários que representam as operações elementares do
  hardware.

  \medskip
  \item \textbf{Unidade central de processamento (CPU)}

  A CPU é a seção ``administrativa'' do computador. Ela coordena e supervisiona o
  funcionamento de todas as demais unidades lógicas: diz à unidade de entrada quando
  ler dados, à ALU quando calcular, à unidade de saída quando enviar resultados, e ao
  armazenamento secundário quando gravar ou recuperar dados.

\end{enumerate}
\end{respostabox}

\newpage

% ============================================================
\section{Exercício 1.6 -- Verdadeiro ou Falso}
% ============================================================

\begin{enunciadoBox}[Enunciado]
Indique se as sentenças a seguir são \textbf{verdadeiras} ou \textbf{falsas}.
Explique sua resposta caso a resposta seja \textit{falsa}.
\begin{enumerate}
  \item Linguagens de máquina geralmente são dependentes da máquina.
  \item Assim como outras linguagens de alto nível, C geralmente é considerada como independente da máquina.
\end{enumerate}
\end{enunciadoBox}

\bigskip

\subsection*{Item (a) -- Linguagens de máquina são dependentes da máquina}

\begin{verdadeiroBox}
\textbf{Sentença: VERDADEIRA}

\medskip
\textbf{Explicação:}

Da Seção~1.6: \textit{``As linguagens de máquina são dependentes da máquina, isto é,
uma linguagem de máquina em particular só pode ser usada em um tipo de computador.''}

Cada família de processadores (Intel x86-64, ARM, RISC-V, MIPS, etc.) possui seu
próprio conjunto de instruções de máquina. Um programa em linguagem de máquina escrito
para um processador Intel não funcionará em um processador ARM, e vice-versa. Isso ocorre
porque a linguagem de máquina é definida pelo projeto físico do hardware do processador
-- pelos circuitos internos que interpretam cada sequência de bits como uma operação
específica.

\medskip
Essa dependência de máquina é exatamente a razão histórica pela qual as linguagens de
alto nível foram criadas: para libertar os programadores da necessidade de reescrever
seus programas inteiramente a cada mudança de hardware.
\end{verdadeiroBox}

\bigskip

\subsection*{Item (b) -- C é independente de máquina}

\begin{verdadeiroBox}
\textbf{Sentença: VERDADEIRA} (com ressalva importante)

\medskip
\textbf{Explicação:}

Da Seção~1.7 e da Dica de Portabilidade~1.1: C é \textit{predominantemente} independente
de hardware. Um programa bem escrito em C pode ser compilado e executado em múltiplas
plataformas (Windows, Linux, macOS, sistemas embarcados) com poucas ou nenhuma
modificação.

\medskip
\textbf{A ressalva essencial -- Dica de Portabilidade~1.2:}
\end{verdadeiroBox}

\begin{portab}
  Embora seja possível escrever programas portáveis em C, existem muitos problemas
  entre compiladores C diferentes e computadores diferentes que podem tornar a
  portabilidade difícil de ser alcançada. Simplesmente escrever programas em C
  \textit{não garante} portabilidade. Com frequência, o programador precisará lidar
  com variações entre computadores.
\end{portab}

Exemplos de situações que podem comprometer a portabilidade em C:
\begin{itemize}
  \item O tamanho em bytes de tipos como \texttt{int} e \texttt{long} varia entre
  plataformas (32 ou 64 bits).
  \item Extensões proprietárias de compiladores (recursos específicos do MSVC ou do GCC
  que não estão no padrão C).
  \item Dependência de comportamento indefinido (\textit{undefined behavior}) do padrão C.
\end{itemize}

\newpage

% ============================================================
\section{Exercício 1.7 -- stdin, stdout e stderr}
% ============================================================

\begin{enunciadoBox}[Enunciado]
Explique o significado dos termos a seguir:
\begin{enumerate}
  \item \texttt{stdin}
  \item \texttt{stdout}
  \item \texttt{stderr}
\end{enumerate}
\end{enunciadoBox}

\bigskip

\begin{respostabox}

Os três termos descrevem os três \textbf{fluxos de dados padrão} (\textit{standard
streams}) disponíveis em todo programa C que é executado em um sistema operacional
baseado em UNIX/Linux (e também em Windows). Eles são introduzidos na Seção~1.14.
Cada fluxo é um canal de comunicação predefinido entre o programa e o sistema operacional.

\medskip

\textbf{a) \texttt{stdin} -- Fluxo de Entrada Padrão (\textit{Standard Input})}

\begin{itemize}
  \item \textbf{O que é:} O canal pelo qual o programa recebe dados de entrada.
  \item \textbf{Dispositivo padrão:} O teclado. Quando seu programa em C chama
  \texttt{scanf()} ou \texttt{getchar()}, ele lê dados do \texttt{stdin}.
  \item \textbf{Redirecionamento:} O \texttt{stdin} pode ser redirecionado para ler de um
  arquivo em vez do teclado. Por exemplo, no terminal Linux:
  \begin{lstlisting}
./meu_programa < dados.txt
  \end{lstlisting}
  Isso faz o programa ler de \texttt{dados.txt} como se o usuário tivesse digitado
  os dados pelo teclado.
  \item \textbf{Origem do nome:} ``stdin'' = \textit{standard input} (entrada padrão).
\end{itemize}

\bigskip

\textbf{b) \texttt{stdout} -- Fluxo de Saída Padrão (\textit{Standard Output})}

\begin{itemize}
  \item \textbf{O que é:} O canal pelo qual o programa envia sua saída normal (resultados,
  mensagens, dados calculados).
  \item \textbf{Dispositivo padrão:} A tela (monitor). Quando seu programa em C chama
  \texttt{printf()}, ele escreve no \texttt{stdout}.
  \item \textbf{Redirecionamento:} O \texttt{stdout} pode ser redirecionado para gravar em
  um arquivo:
  \begin{lstlisting}
./meu_programa > resultado.txt
  \end{lstlisting}
  Isso salva toda a saída do programa no arquivo \texttt{resultado.txt}.
  \item \textbf{Combinação:} É possível redirecionar tanto entrada quanto saída
  simultaneamente:
  \begin{lstlisting}
./meu_programa < dados.txt > resultado.txt
  \end{lstlisting}
\end{itemize}

\bigskip

\textbf{c) \texttt{stderr} -- Fluxo de Erro Padrão (\textit{Standard Error})}

\begin{itemize}
  \item \textbf{O que é:} Um canal separado do \texttt{stdout}, destinado exclusivamente
  a \textit{mensagens de erro e diagnóstico}.
  \item \textbf{Dispositivo padrão:} Também a tela, por padrão.
  \item \textbf{Por que um canal separado?} A grande vantagem de ter \texttt{stderr}
  independente do \texttt{stdout} é que o usuário pode redirecionar a saída normal
  (\texttt{stdout}) para um arquivo, enquanto as mensagens de erro (\texttt{stderr})
  continuam aparecendo na tela -- garantindo que erros nunca passem despercebidos:
  \begin{lstlisting}
./meu_programa > resultado.txt   /* stdout vai para arquivo */
                                 /* stderr continua na tela  */
  \end{lstlisting}
  \item \textbf{Uso em C:} A função \texttt{fprintf(stderr, "Mensagem de erro\n");}
  é usada para escrever no \texttt{stderr}. Isso será estudado em detalhes nos
  capítulos sobre funções e entrada/saída formatada.
\end{itemize}

\medskip
\textbf{Resumo visual:}

\begin{center}
\renewcommand{\arraystretch}{1.4}
\begin{tabular}{>{\ttfamily\bfseries}l l l l}
  \toprule
  \normalfont\textbf{Stream} & \textbf{Tipo} & \textbf{Padrão} & \textbf{Função C} \\
  \midrule
  stdin  & Entrada  & Teclado & \texttt{scanf()}, \texttt{getchar()} \\
  stdout & Saída normal & Tela & \texttt{printf()}, \texttt{putchar()} \\
  stderr & Saída de erros & Tela & \texttt{fprintf(stderr, ...)} \\
  \bottomrule
\end{tabular}
\end{center}

\end{respostabox}

\newpage

% ============================================================
\section{Exercício 1.8 -- Programação Orientada a Objetos}
% ============================================================

\begin{enunciadoBox}[Enunciado]
Por que tanta atenção é dada hoje à programação orientada a objetos?
\end{enunciadoBox}

\bigskip

\begin{respostabox}

Esta é uma questão que vai além de uma resposta técnica simples -- ela toca na história
e na evolução da engenharia de software. Vamos construir a resposta incrementalmente,
a partir do que o Capítulo~1 nos ensina nas Seções~1.9 e~1.13.

\bigskip
\textbf{1. O problema que a POO veio resolver}

Na Seção~1.13, o co-autor Harvey Deitel relata que, já na década de 1960, as organizações
de software enfrentavam graves dificuldades com projetos de grande porte: softwares com
atrasos de entrega, estouro de orçamento e resultados pouco confiáveis. As técnicas de
\textit{programação procedural} (como C, FORTRAN, COBOL) focavam em \textbf{ações
(verbos)}, o que tornava difícil modelar o mundo real -- que é composto de
\textbf{coisas (substantivos)}.

\bigskip
\textbf{2. O que a Programação Orientada a Objetos (POO) oferece}

A POO organiza o software em torno de \textbf{objetos} -- componentes que modelam
entidades do mundo real (uma conta bancária, um produto, um aluno) e encapsulam
tanto seus \textit{dados} quanto seus \textit{comportamentos}. As principais razões
pelas quais a POO recebe tanta atenção são:

\begin{enumerate}
  \item \textbf{Reutilização de software:} Objetos bem projetados (chamados \textit{classes})
  tendem a ser reutilizáveis em projetos futuros. Isso economiza tempo e dinheiro, evitando
  que as equipes ``reinventem a roda'' a cada novo projeto. A Observação sobre Engenharia de
  Software~1.1 do livro destaca que grandes bibliotecas de classes reutilizáveis estão
  disponíveis gratuitamente na Internet.

  \item \textbf{Software mais inteligível e organizado:} Como afirma a Seção~1.13, algumas
  organizações relatam que o principal benefício da POO \textit{não} é a reutilização, mas
  sim que o software produzido é mais \textit{claro, bem organizado e fácil de manter,
  modificar e depurar}. Isso é vital porque até 80\% dos custos de software estão associados
  à \textit{manutenção contínua}, não ao desenvolvimento inicial.

  \item \textbf{Modelagem natural do mundo real:} Vivemos em um mundo de objetos (carros,
  pessoas, produtos, contratos). A POO permite que os programadores modelem o sistema de
  forma que espelhe essa realidade, tornando o design mais intuitivo.

  \item \textbf{Produtividade das equipes:} O uso de uma abordagem modular e orientada
  a objetos torna os grupos de desenvolvimento muito mais produtivos do que seria possível
  com técnicas procedurais.

  \item \textbf{Encapsulamento e proteção de dados:} A POO permite esconder os detalhes
  internos de implementação de um objeto, expondo apenas o que é necessário para o uso
  externo. Isso reduz erros e facilita a evolução do sistema.
\end{enumerate}

\bigskip
\textbf{3. Perspectiva histórica e o papel de C}

C++ (Seção~1.9), desenvolvida por Bjarne Stroustrup no Bell Laboratories, foi a primeira
linguagem a trazer POO com ampla aceitação no mercado, ao adicionar recursos de objetos
à linguagem C. Desde a década de 1990, linguagens orientadas a objetos (C++, Java, C\#,
Python) tornaram-se dominantes -- e a POO passou a ser a metodologia-chave do
desenvolvimento de software.

\medskip
\textbf{Nota importante para este curso:} Você aprenderá primeiro C (Capítulos~1--14),
desenvolvendo uma base sólida em programação estruturada. Nos Capítulos~15--24, você
então estudará C++ e POO -- construindo sobre o alicerce que está construindo agora.

\end{respostabox}

\newpage

% ============================================================
\section{Exercício 1.9 -- Identificando Linguagens de Programação}
% ============================================================

\begin{enunciadoBox}[Enunciado]
Qual linguagem de programação é mais bem descrita por cada uma das seguintes sentenças?
\begin{enumerate}
  \item Desenvolvida pela IBM para aplicações científicas e de engenharia.
  \item Desenvolvida especificamente para aplicações comerciais.
  \item Desenvolvida para o ensino de programação estruturada.
  \item Seu nome homenageia a primeira pessoa a programar um computador no mundo.
  \item Desenvolvida para familiarizar iniciantes com as técnicas de programação.
  \item Desenvolvida especificamente para ajudar os programadores a migrar para .NET.
  \item Conhecida como a linguagem de desenvolvimento do UNIX.
  \item Formada, principalmente, pela inclusão da programação orientada a objetos à linguagem C.
  \item Inicialmente, teve sucesso devido à sua capacidade de criar páginas Web com conteúdo dinâmico.
\end{enumerate}
\end{enunciadoBox}

\bigskip

\begin{respostabox}

\begin{center}
\renewcommand{\arraystretch}{1.6}
\begin{tabular}{>{\bfseries}c p{7.5cm} l}
  \toprule
  \textbf{Item} & \textbf{Descrição} & \textbf{Linguagem} \\
  \midrule
  a & Desenvolvida pela IBM para aplicações científicas e de engenharia. &
    \textcolor{deitelblue}{\bfseries FORTRAN} \\
  b & Desenvolvida especificamente para aplicações comerciais. &
    \textcolor{deitelblue}{\bfseries COBOL} \\
  c & Desenvolvida para o ensino de programação estruturada. &
    \textcolor{deitelblue}{\bfseries Pascal} \\
  d & Seu nome homenageia a primeira pessoa a programar um computador no mundo. &
    \textcolor{deitelblue}{\bfseries Ada} \\
  e & Desenvolvida para familiarizar iniciantes com as técnicas de programação. &
    \textcolor{deitelblue}{\bfseries BASIC} \\
  f & Desenvolvida especificamente para ajudar os programadores a migrar para .NET. &
    \textcolor{deitelblue}{\bfseries Visual C\#} \\
  g & Conhecida como a linguagem de desenvolvimento do UNIX. &
    \textcolor{deitelblue}{\bfseries C} \\
  h & Formada, principalmente, pela inclusão da POO à linguagem C. &
    \textcolor{deitelblue}{\bfseries C++} \\
  i & Teve sucesso pela capacidade de criar páginas Web com conteúdo dinâmico. &
    \textcolor{deitelblue}{\bfseries Java} \\
  \bottomrule
\end{tabular}
\end{center}

\bigskip
\textbf{Explicações detalhadas:}

\textbf{a) FORTRAN (FORmula TRANslator):}
Desenvolvida pela IBM Corporation em meados da década de 1950 (Seção~1.11). Foi
criada para aplicações científicas e de engenharia que exigem cálculos matemáticos
complexos. Ainda é amplamente utilizada em aplicações de engenharia e simulações
científicas.

\medskip
\textbf{b) COBOL (COmmon Business Oriented Language):}
Desenvolvida no final da década de 1950 por fabricantes de computadores, pelo
governo dos EUA e por usuários corporativos (Seção~1.11). É utilizada principalmente
em aplicações comerciais que exigem manipulação precisa de grandes volumes de dados
(como sistemas bancários e de folha de pagamento). Muitos sistemas legados ainda
rodam COBOL hoje.

\medskip
\textbf{c) Pascal:}
Projetada pelo professor Niklaus Wirth em 1971, em homenagem ao matemático Blaise
Pascal (Seção~1.11). Foi criada especificamente para o ensino de programação
estruturada e tornou-se a linguagem preferida em ambientes acadêmicos por décadas.

\medskip
\textbf{d) Ada:}
Desenvolvida com o patrocínio do Departamento de Defesa dos EUA na década de 1970
(Seção~1.11). Recebeu o nome de \textbf{Lady Ada Lovelace}, filha do poeta Lord
Byron, reconhecida historicamente como a primeira programadora do mundo (século XIX).
Ada suporta \textit{multitarefa} nativamente.

\medskip
\textbf{e) BASIC (Beginner's All-purpose Symbolic Instruction Code):}
Desenvolvida em meados da década de 1960 no Dartmouth College (Seção~1.12), com o
objetivo de familiarizar iniciantes com as técnicas de programação. Foi a linguagem
de entrada de inúmeros programadores nas décadas de 1970 e 1980.

\medskip
\textbf{f) Visual C\#:}
Uma nova linguagem criada pela Microsoft exclusivamente para a plataforma .NET
(Seção~1.12), baseada em C++ e Java. Foi desenvolvida especificamente para ajudar
os programadores a criar aplicações para o ambiente .NET, facilitando a migração
de desenvolvedores C++ e Java para a plataforma Microsoft.

\medskip
\textbf{g) C:}
Como já estudamos extensamente no Capítulo~1 (Seção~1.7), C ficou mundialmente
conhecida como a linguagem de desenvolvimento do sistema operacional UNIX. Criada por
Dennis Ritchie no Bell Laboratories em 1972, C e UNIX cresceram juntos e definiram
uma era na computação.

\medskip
\textbf{h) C++:}
Desenvolvida por Bjarne Stroustrup no Bell Laboratories (Seção~1.9). C++ foi
formada adicionando-se recursos de programação orientada a objetos à linguagem C
-- especificamente, absorvendo os conceitos de objetos da linguagem de simulação
Simula~67.

\medskip
\textbf{i) Java:}
Criada pela Sun Microsystems com base em C++ (Seção~1.10). Embora inicialmente
desenvolvida para aparelhos eletrônicos, Java ganhou enorme popularidade quando
a World Wide Web explodiu em 1993, pois permitia criar \textit{applets} com conteúdo
dinâmico e interativo nas páginas Web -- uma capacidade revolucionária na época.

\end{respostabox}

\newpage

% ============================================================
\section*{Resumo das Respostas -- Exercícios 1.3 a 1.9}
% ============================================================

\begin{tcolorbox}[colback=deitellightblue, colframe=deitelblue,
  title={\bfseries\large Gabarito Rápido}, breakable]

\textbf{1.3 -- Hardware ou Software:}
\begin{itemize}[noitemsep]
  \item[a)] CPU $\to$ \textbf{hardware} \quad
  b) compilador C++ $\to$ \textbf{software} \quad
  c) ALU $\to$ \textbf{hardware}
  \item[d)] pré-processador C++ $\to$ \textbf{software} \quad
  e) unidade de entrada $\to$ \textbf{hardware} \quad
  f) programa editor $\to$ \textbf{software}
\end{itemize}

\medskip
\textbf{1.5 -- Unidades do Computador:}
\begin{itemize}[noitemsep]
  \item[a)] unidade de entrada \quad b) programação \quad c) linguagem simbólica
  \item[d)] unidade de saída \quad e) memória + armazenamento secundário
  \item[f)] ALU \quad g) ALU \quad h) linguagem de alto nível
  \item[i)] linguagem de máquina \quad j) CPU
\end{itemize}

\medskip
\textbf{1.6 -- Verdadeiro/Falso:}
\begin{itemize}[noitemsep]
  \item[a)] \textbf{Verdadeiro} -- linguagens de máquina são dependentes da máquina.
  \item[b)] \textbf{Verdadeiro} (com ressalva) -- C é predominantemente independente,
  mas portabilidade perfeita exige cuidado.
\end{itemize}

\medskip
\textbf{1.7 -- Streams:}
\begin{itemize}[noitemsep]
  \item \texttt{stdin} = fluxo de entrada padrão (teclado)
  \item \texttt{stdout} = fluxo de saída padrão (tela)
  \item \texttt{stderr} = fluxo de erro padrão (tela, separado do stdout)
\end{itemize}

\medskip
\textbf{1.9 -- Linguagens:}
\begin{itemize}[noitemsep]
  \item[a)] FORTRAN \quad b) COBOL \quad c) Pascal \quad d) Ada \quad e) BASIC
  \item[f)] Visual C\# \quad g) C \quad h) C++ \quad i) Java
\end{itemize}

\end{tcolorbox}

\bigskip

\begin{boaprática}
  Você completou os Exercícios do Capítulo~1! Os conceitos aqui revisados --
  organização do computador, tipos de linguagens, ambiente de desenvolvimento,
  streams padrão -- formam a base conceitual sobre a qual todo o restante do
  livro é construído. No Capítulo~2, você começará a escrever seus primeiros
  programas em C, aplicando esses conceitos de forma prática e concreta.
\end{boaprática}

\end{document}
